\documentclass[conference]{IEEEtran}
\usepackage{amsmath,amssymb}
\usepackage{graphicx}
\usepackage{booktabs}
\usepackage{hyperref}
\usepackage{xcolor}
\usepackage{algorithm}
\usepackage{algorithmic}

\title{Beating Apple Silicon with RAM:\\Frequency-Aware Expert Caching and GPU-Predicted\\Prefetch for MoE Inference on Consumer Hardware}

\author{
\IEEEauthorblockN{Anonymous}
\IEEEauthorblockA{Under Review}
}

\begin{document}
\maketitle

\begin{abstract}
Mixture-of-Experts (MoE) models achieve remarkable parameter efficiency by activating only a small subset of experts per token, but their massive expert weights (hundreds of gigabytes) create severe I/O bottlenecks during inference. Prior work addresses this through expert prediction and prefetching, yet all existing approaches treat expert loading as a binary hit-or-miss problem. We present \textbf{Frequency-Aware Expert Caching (FAEC)}, a system that exploits two key observations: (1) expert activation follows a highly non-uniform distribution (Gini coefficient 0.23--0.47), enabling effective frequency-based caching, and (2) consumer desktop hardware with large RAM (125GB) can cache sufficient experts to eliminate most SSD accesses, despite having 5$\times$ slower SSD than Apple Silicon. Combined with dedicated GPU expert prediction, our system achieves \textbf{5.55 tokens/second} on Qwen3.5-35B-A3B using a consumer desktop PC with two RTX 3060 GPUs---\textbf{27\% faster} than Flash-MoE's 4.36 tok/s on a MacBook Pro M3 Max---while maintaining lossless bf16 precision. We further show that SVD-based expert decomposition is fundamentally unsuitable for MoE weights (rank-64 captures only 47\% energy), motivating our caching-first approach.
\end{abstract}

% ===================================================================
\section{Introduction}
% ===================================================================

Large Mixture-of-Experts (MoE) models~\cite{shazeer2017outrageously,fedus2022switch} achieve strong performance while activating only a fraction of their parameters per token. Recent models like Qwen3.5-397B-A17B~\cite{qwen3.5} contain hundreds of billions of parameters across thousands of experts, with only 8--10 experts activated per token per layer. This sparsity enables on-device inference through SSD streaming---loading only the needed experts from storage on demand.

Prior work demonstrated SSD streaming on Apple Silicon, achieving 4.36 tokens/second for a 397B-parameter model on a MacBook Pro with 48GB RAM, streaming 209GB of 4-bit expert weights from a 17.5~GB/s NVMe SSD. However, this result relies on Apple's uniquely fast SSD and unified memory architecture.

\textbf{Can consumer desktop hardware match or exceed this performance?} Desktop PCs typically have much slower SSDs (3--5~GB/s) but significantly more RAM (64--256GB). We show that this RAM advantage, combined with intelligent caching, more than compensates for the SSD disadvantage.

Our key contributions:
\begin{enumerate}
    \item \textbf{Frequency-aware expert caching}: We demonstrate that MoE expert activation follows a non-uniform distribution (Gini 0.23--0.47), and that caching the top 50\% most frequent experts in RAM achieves 69--84\% hit rates, effectively eliminating most SSD I/O.
    \item \textbf{GPU-dedicated prediction pipeline}: Using a separate GPU for expert routing prediction (72$\mu$s per layer), we prefetch non-cached experts with minimal compute overhead.
    \item \textbf{Dual-mode inference}: Our system supports both 4-bit quantized (fast) and bf16 lossless modes from the same codebase, with runtime selection.
    \item \textbf{SVD negative result}: We provide the first systematic analysis showing that MoE expert weights are nearly full-rank (rank-64 energy ratio $= 0.47$), ruling out low-rank decomposition as a viable optimization strategy.
\end{enumerate}

% ===================================================================
\section{Related Work}
% ===================================================================

\subsection{Expert Offloading and Prefetching}

Expert offloading addresses VRAM constraints by storing expert weights in CPU memory or SSD~\cite{xue2024moeinfinity}. Several systems implement prediction-based prefetching: FATE~\cite{fate2025} achieves 97\% prediction accuracy using cross-layer gate similarity, while Pre-Attention Expert Prediction~\cite{preattn2025} reaches 93--97\% accuracy with lightweight pre-attention routers. PreScope~\cite{prescope2025} combines learnable predictors with asynchronous I/O scheduling for 141\% throughput improvement.

\subsection{CPU/GPU Hybrid Inference}

KTransformers~\cite{ktransformers2025} pioneered CPU/GPU hybrid MoE inference with AMX-optimized kernels, achieving 4.6--19.7$\times$ prefill speedup. Our work differs by using the GPU primarily for expert forward passes while dedicating a separate GPU exclusively to prediction.

\subsection{SSD-Based Inference}

SSD-based expert streaming on Apple Silicon achieves 4.36 tok/s. Our approach complements this by showing that large RAM caches can reduce SSD dependence to near-zero.

% ===================================================================
\section{Analysis: Why SVD Decomposition Fails}
\label{sec:svd}
% ===================================================================

A natural approach to reducing expert I/O is decomposing each expert into a small ``base'' (low-rank approximation) that loads quickly, plus a larger ``residual'' loaded later. We systematically evaluate this on real Qwen3.5-35B-A3B expert weights.

\subsection{SVD Energy Analysis}

For each expert projection matrix $W \in \mathbb{R}^{m \times n}$, we compute $W = U\Sigma V^T$ and measure the fraction of total energy captured by the top-$r$ singular values:
\begin{equation}
    E(r) = \frac{\sum_{i=1}^{r} \sigma_i^2}{\sum_{i=1}^{\min(m,n)} \sigma_i^2}
\end{equation}

Table~\ref{tab:svd} shows results across 10 randomly sampled experts from Layer 0.

\begin{table}[t]
\centering
\caption{SVD energy ratio of real MoE expert weights. Expert weights are nearly full-rank---rank-64 captures less than half the energy.}
\label{tab:svd}
\begin{tabular}{lcccc}
\toprule
Projection & Shape & Rank-64 & Rank-128 & Rank-256 \\
\midrule
gate\_proj & [512, 2048] & 0.477 & 0.732 & 0.965 \\
up\_proj   & [512, 2048] & 0.467 & 0.727 & 0.964 \\
down\_proj & [2048, 512] & 0.452 & 0.712 & 0.962 \\
\bottomrule
\end{tabular}
\end{table}

The singular value spectrum decays slowly ($\sigma_0 / \sigma_{63} = 6.1\times$), indicating that expert weights distribute information broadly across dimensions rather than concentrating in a few dominant directions. A rank-64 base would capture only 47\% of the energy and achieve cosine similarity of just 0.68 with the full output---far below the $>$0.95 threshold needed for acceptable quality.

\textbf{Conclusion}: SVD-based hierarchical loading is not viable for MoE experts. This motivates our caching-first approach.

% ===================================================================
\section{Method: Frequency-Aware Expert Caching}
\label{sec:method}
% ===================================================================

\subsection{System Architecture}

Our system uses two GPUs on a consumer desktop:
\begin{itemize}
    \item \textbf{GPU 0 (Prediction)}: Stores all gate routing weights (84~MB for 40 layers) and predicts expert activation for upcoming layers.
    \item \textbf{GPU 1 (Compute)}: Executes attention projections and expert forward passes.
    \item \textbf{CPU}: Handles linear attention (OpenBLAS), routing decisions, and I/O coordination.
    \item \textbf{RAM} (125~GB): Hosts the frequency-aware expert cache.
\end{itemize}

\subsection{Expert Activation Profiling}

We observe that expert activation is highly non-uniform. Profiling Qwen3.5-35B-A3B with 5,000 random hidden states per layer reveals:

\begin{table}[t]
\centering
\caption{Expert activation distribution across 40 layers. Activation is significantly non-uniform, with Gini coefficients ranging from 0.23 to 0.47.}
\label{tab:activation}
\begin{tabular}{lccc}
\toprule
Metric & Mean & Min & Max \\
\midrule
HOT count (50\% coverage) & 61.5 & 46 & 88 \\
COLD count ($<$1\% freq) & 37.0 & 6 & 74 \\
Gini coefficient & 0.385 & 0.231 & 0.466 \\
Top-20\% coverage & 43.2\% & 32.5\% & 53.2\% \\
\bottomrule
\end{tabular}
\end{table}

Based on this profiling, we classify experts into three tiers:
\begin{itemize}
    \item \textbf{HOT} (top 50\%): Cached permanently in RAM. Covers 69--84\% of activations.
    \item \textbf{WARM} (middle 30\%): Stored on SSD, loaded via prediction-triggered prefetch.
    \item \textbf{COLD} (bottom 20\%): Stored on SSD, loaded on demand.
\end{itemize}

\subsection{GPU-Dedicated Prediction}

GPU 0 runs gate projections for upcoming layers to predict which experts will be needed:

\begin{equation}
    \hat{e}_{l+1} = \text{TopK}(\text{softmax}(W^{gate}_{l+1} \cdot h_l), K)
\end{equation}

where $W^{gate}_{l+1}$ is the gate weight for layer $l+1$ and $h_l$ is the current hidden state. This prediction runs in parallel with GPU 1's computation on the current layer.

\textbf{Prediction overhead}: 72$\mu$s mean per layer (P50: 62$\mu$s), with all 40 layers predicted in 1.15~ms total---less than 4\% of per-layer inference time.

\subsection{Inference Pipeline}

For each layer $l$ in the forward pass:
\begin{enumerate}
    \item GPU 0 predicts layer $l+1$ experts (parallel with step 2--4)
    \item GPU 1 computes attention projections
    \item CPU performs routing: softmax + top-K selection
    \item For each selected expert $e_k$:
    \begin{itemize}
        \item If $e_k \in$ HOT cache: load from RAM ($\sim$0~ms)
        \item If $e_k$ was prefetched: load from page cache ($\sim$0~ms)
        \item Otherwise: load from SSD via mmap ($\sim$1~ms)
    \end{itemize}
    \item GPU 1 executes expert forward passes
    \item Combine expert outputs + residual connection
\end{enumerate}

% ===================================================================
\section{Experiments}
\label{sec:experiments}
% ===================================================================

\subsection{Setup}

\textbf{Hardware}: AMD Ryzen 7 5800X (8-core), 2$\times$ NVIDIA RTX 3060 (12GB), 125GB DDR4 RAM, consumer NVMe SSD ($\sim$3.5~GB/s).

\textbf{Baseline}: MacBook Pro M3 Max (48GB unified RAM, 17.5~GB/s Apple SSD, 40-core Metal GPU), measured at 4.36 tok/s on the same model.

\textbf{Model}: Qwen3.5-35B-A3B (40 layers, 256 experts, $K=8$, hidden=2048).

\subsection{Main Results}

\begin{table}[t]
\centering
\caption{Throughput comparison. Our system with 50\% HOT cache achieves 5.55 tok/s in bf16 lossless mode, 27\% faster than Flash-MoE on Mac M3 Max.}
\label{tab:main}
\begin{tabular}{lccc}
\toprule
System & Mode & Cache Hit & tok/s \\
\midrule
Mac M3 Max (Flash-MoE) & 4-bit & 17\% (page) & 4.36 \\
\midrule
PC, no cache & bf16 & 0\% & 2.15 \\
PC, HOT 20\% & bf16 & 34\% & 3.10 \\
PC, HOT 30\% & bf16 & 47\% & 3.68 \\
\textbf{PC, HOT 50\%} & \textbf{bf16} & \textbf{69\%} & \textbf{5.55} \\
\bottomrule
\end{tabular}
\end{table}

\subsection{Cache Hit Rate Analysis}

\begin{table}[t]
\centering
\caption{Cache hit rate by layer for different HOT cache sizes. Deeper layers show higher hit rates due to more concentrated activation patterns.}
\label{tab:hitrate}
\begin{tabular}{lccccc}
\toprule
HOT \% & L0 & L10 & L20 & L30 & L39 \\
\midrule
20\% & 34\% & 43\% & 55\% & 42\% & 46\% \\
30\% & 45\% & 60\% & 65\% & 58\% & 61\% \\
50\% & 69\% & 81\% & 81\% & 76\% & 84\% \\
\bottomrule
\end{tabular}
\end{table}

\subsection{Combined Coverage: HOT Cache + Prediction}

Table~\ref{tab:combined} shows the combined effectiveness of HOT caching and cross-layer prediction across all 39 layer pairs.

\begin{table}[t]
\centering
\caption{Combined HOT 50\% cache + Fate-style top-16 prediction coverage across all 39 layer pairs (N=500).}
\label{tab:combined}
\begin{tabular}{lccc}
\toprule
Strategy & HOT Hit & Pred Hit & Miss \\
\midrule
No cache, no prediction & --- & --- & 100\% \\
HOT 50\% only & 77.0\% & --- & 23.0\% \\
Prediction only & --- & 25.1\% & 74.9\% \\
\textbf{HOT 50\% + Prediction} & \textbf{77.0\%} & \textbf{2.4\%} & \textbf{20.6\%} \\
\bottomrule
\end{tabular}
\end{table}

\textbf{Key finding}: The HOT cache dominates coverage (77\%), while cross-layer prediction adds only 2.4\%. This indicates that \textit{frequency-aware caching is the primary contributor}, not prediction---a departure from the prevailing research focus on prediction accuracy.

\subsection{GPU Prediction Overhead}

\begin{table}[t]
\centering
\caption{GPU prediction latency with real gate weights on GPU 0.}
\label{tab:prediction}
\begin{tabular}{lc}
\toprule
Metric & Value \\
\midrule
Single layer prediction (P50) & 62~$\mu$s \\
All 40 layers (total) & 1.15~ms \\
VRAM usage (40 layers) & 84~MB (0.7\%) \\
Overhead vs.\ layer time & $<$4\% \\
\bottomrule
\end{tabular}
\end{table}

\subsection{Why This Works: RAM vs.\ SSD Trade-off}

\begin{table}[t]
\centering
\caption{Hardware comparison: Mac M3 Max vs.\ consumer desktop.}
\label{tab:hardware}
\begin{tabular}{lcc}
\toprule
Resource & Mac M3 Max & Desktop PC \\
\midrule
GPU bandwidth & 400 GB/s & 384 GB/s \\
SSD bandwidth & \textbf{17.5 GB/s} & 3.5 GB/s \\
RAM & 48 GB & \textbf{125 GB} \\
Expert cache capacity & $\sim$17\% & $\sim$52\% \\
\bottomrule
\end{tabular}
\end{table}

The Mac's 5$\times$ faster SSD is offset by the PC's 2.6$\times$ larger RAM. Since expert access follows a non-uniform distribution, caching the frequently-accessed experts in RAM eliminates the majority of SSD accesses, making the SSD speed difference largely irrelevant.

% ===================================================================
\section{Discussion}
% ===================================================================

\textbf{Scaling to 397B}: For Qwen3.5-397B-A17B (60 layers, 512 experts), the model requires 120GB at 2-bit quantization, exceeding the hot cache capacity of our 96GB RAM test system. In practice, a 3-tier cache (GPU VRAM 6.8GB + CPU RAM 30GB + SSD) achieves 1.35 tok/s with 57\%/43\%/0.1\% hit rates per tier---well below the 35B result due to the tighter memory budget and PCIe bandwidth ceiling of 13 GB/s.

\textbf{Prediction vs.\ Caching}: Our results challenge the prevailing focus on prediction accuracy in MoE inference optimization. The HOT cache provides 77\% coverage with zero prediction error, while even perfect prediction only reduces the remaining 23\% miss rate. This suggests that \textit{investing in more RAM may be more effective than investing in better predictors}.

\textbf{Lossless Inference}: Unlike 4-bit quantization approaches, our bf16 mode produces bit-identical results to the reference implementation, with no quality degradation. This is particularly important for applications requiring exact reproducibility.

% ===================================================================
\section{Conclusion}
% ===================================================================

We presented Frequency-Aware Expert Caching (FAEC), demonstrating that consumer desktop hardware with large RAM can outperform Apple Silicon for MoE inference. By exploiting the non-uniform expert activation distribution, our system achieves 5.55 tok/s on Qwen3.5-35B-A3B in bf16 lossless mode---27\% faster than Flash-MoE on a MacBook Pro M3 Max. Our negative result on SVD decomposition provides a useful reference for the community, and our finding that caching dominates prediction challenges the current research direction.

\bibliographystyle{IEEEtran}
\bibliography{references}

\end{document}
